% Unit 3 self-study -- "I do / You do" ladder for CED 3.1-3.13.
% Fourteen ladders, fifty-six questions: the heaviest unit on the exam
% (18-22%) gets the longest self-study.
%
% Disjoint from u03-frq, which uses the Zn/HCl gas-collection question,
% the CH4/CH3F/CH3OH boiling-point ranking, the real-gas two-regime
% question and an A=0.435 Beer-Lambert item. This document uses NaHCO3
% decomposition, a solids-classification ladder, effusion ratios and
% A=0.620 instead.
\documentclass{apchem-worksheet}

\apcunit{3}
\apcunittitle{Properties of Substances and Mixtures}
\apcdoctitle{Self-Study \textbullet\ CED 3.1--3.13, I~do / You~do}
\apcsubtitle{Fourteen ladders \textbullet\ four YOUR TURN questions each \textbullet\ the heaviest unit on the exam}

\begin{document}
\apctitleblock

\begin{apcnote}[How to use these notes --- read this first]
Each skill is a \textbf{ladder}: a worked example, then \textbf{four}
YOUR TURN questions --- same skill, new numbers, no help. Work all four
before checking anything.

\textbf{This unit is worth \SI{18}{\percent}--\SI{22}{\percent} of the AP
exam}, roughly one question in five, and it is assembled from seven
different places in the textbook. That is why it is the longest self-study
in the course. Do not try to do all fourteen ladders in one sitting.

\textbf{Two habits that decide this unit.} Temperature in any gas
calculation is \textbf{always kelvin}. And when you explain why one
substance boils higher, melts higher or dissolves better than another,
name the \textbf{intermolecular force} and say why it is stronger --- never
just assert that it is.

\textbf{One scope note.} Phase diagrams and colligative properties
(freezing-point depression, boiling-point elevation, osmotic pressure) are
\textbf{excluded} by the CED, and molality and percent-by-mass
calculations with them. Concentration here means \textbf{molarity}.
\end{apcnote}

% =====================================================================
\section*{\sffamily Ladder 1 \textbullet\ Identifying intermolecular forces}
\ced{3.1}

Three forces, in increasing strength. \textbf{London dispersion} acts in
every substance. \textbf{Dipole--dipole} adds for polar molecules.
\textbf{Hydrogen bonding} requires H bonded \emph{directly} to N, O or F.
A fourth, \textbf{ion--dipole}, appears when an ion is dissolved.

\begin{workedexample}[I do: naming the strongest force present]
Identify the strongest intermolecular force in \ce{CH4}, \ce{CH2Cl2},
\ce{NH3} and a solution of \ce{NaCl} in water.

\textbf{\ce{CH4}} --- nonpolar, symmetric. Only \textbf{London
dispersion}.

\textbf{\ce{CH2Cl2}} --- polar (the shape does not cancel the \ce{C-Cl}
dipoles), but its hydrogens are on carbon. \textbf{Dipole--dipole}.

\textbf{\ce{NH3}} --- hydrogen bonded directly to nitrogen.
\textbf{Hydrogen bonding}.

\textbf{\ce{NaCl}(aq)} --- ions surrounded by polar water molecules.
\textbf{Ion--dipole}, the strongest of the set.

\textbf{The trap:} calling any molecule containing hydrogen a
hydrogen-bonder. The hydrogen must be attached to N, O or F --- carbon does
not count, however many hydrogens there are.
\end{workedexample}

\begin{yourturn}[YOUR TURN 1 --- four questions]
\begin{enumerate}[leftmargin=*,itemsep=0.5em]
  \item Strongest force in \ce{CO2}: \workspace[1.4cm]
  \item Strongest force in \ce{CH3OH}: \workspace[1.4cm]
  \item Strongest force in \ce{HCl}: \workspace[1.4cm]
  \item \ce{CH3F} contains fluorine but cannot hydrogen bond. Explain.
        \workspace[1.8cm]
\end{enumerate}
\selfcheck{(a) London dispersion \quad (b) hydrogen bonding \quad
(c) dipole--dipole \quad (d) its H atoms are on carbon, not on F}
\keyonly{\begin{apcanswer}(a) \textbf{London dispersion} --- \ce{CO2} is
linear and nonpolar, so its bond dipoles cancel.
(b) \textbf{Hydrogen bonding} --- the \ce{-OH} group puts a hydrogen
directly on oxygen.
(c) \textbf{Dipole--dipole} --- polar, but the hydrogen is on chlorine's
partner, not on N, O or F.
(d) Hydrogen bonding requires a hydrogen bonded \emph{directly} to N, O or
F. In \ce{CH3F} the fluorine is bonded to \emph{carbon}, and all three
hydrogens are on carbon too. There is no \ce{H-F} bond, so no hydrogen
bond can form --- the molecule manages only dipole--dipole
attraction.\end{apcanswer}}
\end{yourturn}

% =====================================================================
\section*{\sffamily Ladder 2 \textbullet\ Ranking boiling points}
\ced{3.1} \ced{3.3}

Boiling means separating molecules completely, so the boiling point tracks
the strength of the intermolecular forces. Compare force \emph{type}
first; only if the types match does size (electron count) decide.

\begin{workedexample}[I do: two comparisons, two different reasons]
Rank \ce{Ar} ($-186$), \ce{HCl} ($-85$) and \ce{H2O} ($+100$
\si{\celsius}), then explain why \ce{Br2} boils higher than \ce{Cl2}.

\textbf{The three-way ranking is by force type.} \ce{Ar} has dispersion
only; \ce{HCl} adds dipole--dipole; \ce{H2O} hydrogen bonds. Stronger
forces need more energy to overcome, so the order follows directly.

\textbf{\ce{Br2} against \ce{Cl2} is a different comparison.} Both are
nonpolar with dispersion only, so type cannot decide. \ce{Br2} has far
more electrons, so its electron cloud is more \textbf{polarizable} --- it
distorts more easily, giving larger instantaneous dipoles and stronger
dispersion forces. Hence the higher boiling point.

\textbf{The method:} ask \emph{is the force type the same?} If not, type
decides. If so, count electrons.
\end{workedexample}

\begin{yourturn}[YOUR TURN 2 --- four questions]
\begin{enumerate}[leftmargin=*,itemsep=0.5em]
  \item Which boils higher, \ce{CH4} or \ce{C4H10}? Why?
        \workspace[1.6cm]
  \item Which boils higher, \ce{H2S} or \ce{H2O}? Why?
        \workspace[1.6cm]
  \item Which boils higher, \ce{CH3OCH3} or \ce{CH3CH2OH}? Both are
        \ce{C2H6O}. \workspace[1.8cm]
  \item \ce{HF} has stronger hydrogen bonds than \ce{H2O}, yet water boils
        higher. Explain. \workspace[2.0cm]
\end{enumerate}
\selfcheck{(a) \ce{C4H10} --- more electrons \quad (b) \ce{H2O} ---
hydrogen bonding \quad (c) ethanol --- its H is on O \quad (d) water forms
more hydrogen bonds per molecule}
\keyonly{\begin{apcanswer}(a) \textbf{\ce{C4H10}}. Both are nonpolar with
dispersion only, so electron count decides; butane has far more electrons
and a more polarizable cloud.
(b) \textbf{\ce{H2O}}. Sulfur is too large and too weakly electronegative
for \ce{H2S} to hydrogen bond, so it manages only dipole--dipole; water
hydrogen bonds. Force type beats the fact that \ce{H2S} has more
electrons.
(c) \textbf{Ethanol}. Same formula and the same number of electrons, so
dispersion is comparable --- but ethanol's hydrogen sits on the
\emph{oxygen} and can hydrogen bond, while both of the ether's hydrogens
are on carbon. The whole difference is hydrogen bonding.
(d) Because the \emph{number} of hydrogen bonds matters more here than the
strength of any one. Water has two hydrogens and two lone pairs, so it can
form up to four hydrogen bonds per molecule; \ce{HF} has one hydrogen and
three lone pairs, so it is limited to about two.\end{apcanswer}}
\end{yourturn}

% =====================================================================
\section*{\sffamily Ladder 3 \textbullet\ Classifying solids}
\ced{3.2}

Four types, each with properties that follow from what is holding the
particles together.

\begin{workedexample}[I do: classify and predict]
Classify \ce{SiO2}, \ce{NaCl}, \ce{Cu} and solid \ce{CO2}, and give one
property of each that follows.

\textbf{\ce{SiO2} --- covalent network.} A continuous lattice of strong
covalent bonds. Melting means breaking covalent bonds, so the melting
point is very high and it is hard and non-conducting.

\textbf{\ce{NaCl} --- ionic.} Strong electrostatic attractions between
ions: high-melting and brittle, conducting only when molten or dissolved.

\textbf{\ce{Cu} --- metallic.} Cations in a delocalized electron sea:
conducts as a solid, and is malleable.

\textbf{Solid \ce{CO2} --- molecular.} Only weak dispersion forces between
whole \ce{CO2} molecules. Very low melting point --- it sublimes at
$-78$~\si{\celsius} --- and does not conduct.

\textbf{The distinction that matters most:} in a molecular solid the
covalent bonds \emph{within} each molecule are strong, but melting only
has to overcome the weak forces \emph{between} molecules. That is why dry
ice sublimes so easily despite containing strong \ce{C=O} bonds.
\end{workedexample}

\begin{yourturn}[YOUR TURN 3 --- four questions]
\begin{enumerate}[leftmargin=*,itemsep=0.5em]
  \item Classify diamond and give its expected melting point behaviour.
        \workspace[1.6cm]
  \item Classify solid \ce{I2}: \workspace[1.4cm]
  \item A solid is hard, high-melting, and conducts only when molten.
        Which type? \workspace[1.6cm]
  \item Explain why graphite conducts electricity but diamond does not,
        though both are carbon. \workspace[2.0cm]
\end{enumerate}
\selfcheck{(a) covalent network, very high \quad (b) molecular \quad
(c) ionic \quad (d) graphite has delocalized electrons; diamond's are all
in sigma bonds}
\keyonly{\begin{apcanswer}(a) \textbf{Covalent network}. Melting requires
breaking a continuous lattice of strong \ce{C-C} covalent bonds, so the
melting point is extremely high (above \SI{3500}{\celsius}).
(b) \textbf{Molecular} --- discrete \ce{I2} molecules held by dispersion
forces only, so it is soft and low-melting.
(c) \textbf{Ionic}. Hardness and a high melting point indicate strong
attractions; conducting only when molten indicates charged particles that
are fixed in the solid and mobile in the melt.
(d) In \textbf{diamond} every carbon forms four single bonds, so all four
valence electrons are localized in sigma bonds and none is free to move.
In \textbf{graphite} each carbon bonds to only three others, leaving one
electron per atom \textbf{delocalized} over the layer. Those mobile
electrons carry current along the sheets.\end{apcanswer}}
\end{yourturn}

% =====================================================================
\section*{\sffamily Ladder 4 \textbullet\ Phase changes and heating curves}
\ced{3.3}

On a heating curve, sloped sections are temperature rising (kinetic energy
increasing) and flat sections are phase changes (potential energy
increasing while temperature holds constant).

\begin{workedexample}[I do: reading the flat sections]
A heating curve for water runs from ice below \SI{0}{\celsius} to steam
above \SI{100}{\celsius}. Explain the two flat regions, and why the second
is longer.

\textbf{The flats are phase changes.} At \SI{0}{\celsius} the ice melts;
at \SI{100}{\celsius} the water boils. Throughout each, energy is being
supplied but the temperature does not move.

\textbf{Why the temperature holds:} the energy goes into overcoming
intermolecular attractions --- raising \emph{potential} energy --- not into
speeding the molecules up. Temperature measures average \emph{kinetic}
energy, which is unchanged.

\textbf{Why the boiling flat is longer:} melting only loosens the lattice
--- molecules leave fixed positions but stay in contact, and most hydrogen
bonds survive. Vaporizing must separate them completely, overcoming
essentially all the attractions. For water that is
\SI{2260}{\joule\per\gram} against \SI{334}{\joule\per\gram}, nearly seven
times as much.
\end{workedexample}

\begin{yourturn}[YOUR TURN 4 --- four questions]
\begin{enumerate}[leftmargin=*,itemsep=0.5em]
  \item On a heating curve, what is happening on a \emph{sloped} section?
        \workspace[1.6cm]
  \item Why does temperature stay constant during melting?
        \workspace[1.8cm]
  \item Which requires more energy per gram for water, melting or
        vaporizing, and why? \workspace[1.8cm]
  \item Steam at \SI{100}{\celsius} causes worse burns than water at
        \SI{100}{\celsius}. Explain. \workspace[2.0cm]
\end{enumerate}
\selfcheck{(a) temperature rising, kinetic energy increasing \quad
(b) energy goes to potential, not kinetic \quad (c) vaporizing ---
separates molecules completely \quad (d) it releases its heat of
condensation as well}
\keyonly{\begin{apcanswer}(a) The substance is in a single phase and its
temperature is rising: the energy supplied increases the average
\emph{kinetic} energy of the particles.
(b) The energy supplied is used to overcome the intermolecular
attractions holding the lattice, raising the \textbf{potential} energy.
Since temperature reflects average kinetic energy, and that is not
changing, the temperature holds steady.
(c) \textbf{Vaporizing}, by a factor of nearly seven
(\SI{2260}{} against \SI{334}{\joule\per\gram}). Melting only loosens the
structure; vaporizing separates the molecules completely.
(d) Steam must first \emph{condense} on the skin, releasing its full heat
of vaporization --- about \SI{2260}{\joule\per\gram} --- before it even
begins to cool. Liquid water at the same temperature delivers only the
energy from cooling. The condensation step is the extra
damage.\end{apcanswer}}
\end{yourturn}

% =====================================================================
\section*{\sffamily Ladder 5 \textbullet\ The ideal gas law}
\ced{3.4}

$PV = nRT$, with $R = \SI{0.08206}{\litre\atm\per\mole\per\kelvin}$. That
choice of $R$ fixes every other unit: pressure in atm, volume in litres,
temperature in \textbf{kelvin}.

\begin{workedexample}[I do: solving for volume]
What volume does \SI{0.250}{\mole} of gas occupy at \SI{2.00}{\atm} and
\SI{300}{\kelvin}?

\[ V = \frac{nRT}{P} = \frac{(0.250)(0.08206)(300)}{2.00}
   = \mathbf{3.08~L} \]

\textbf{Sense check:} one mole at \SI{1}{\atm} and \SI{273}{\kelvin} is
\SI{22.4}{\litre}. A quarter of a mole at double the pressure should be
roughly $22.4/4/2 \approx 2.8$~L, and a slightly higher temperature nudges
it up. \num{3.08} is right.

\textbf{The error that costs whole questions:} leaving the temperature in
\si{\celsius}. At \SI{27}{\celsius} you must use \SI{300}{\kelvin}, and
using 27 gives an answer wrong by a factor of eleven.
\end{workedexample}

\begin{yourturn}[YOUR TURN 5 --- four questions]
\begin{enumerate}[leftmargin=*,itemsep=0.5em]
  \item \SI{1.50}{\mole} in \SI{5.00}{\litre} at \SI{350}{\kelvin}.
        Pressure? \workspace[1.6cm]
  \item \SI{2.00}{\atm}, \SI{4.00}{\litre}, \SI{0.300}{\mole}.
        Temperature? \workspace[1.6cm]
  \item \SI{0.500}{\gram} of helium in \SI{2.00}{\litre} at
        \SI{298}{\kelvin}. Pressure? \workspace[1.8cm]
  \item A rigid container is heated from \SI{300}{\kelvin} to
        \SI{600}{\kelvin}. What happens to the pressure, and why?
        \workspace[1.8cm]
\end{enumerate}
\selfcheck{(a) \SI{8.62}{\atm} \quad (b) \SI{325}{\kelvin} \quad
(c) \SI{1.53}{\atm} \quad (d) doubles --- $P \propto T$ at fixed $n$, $V$}
\keyonly{\begin{apcanswer}(a) $P = nRT/V = (1.50)(0.08206)(350)/5.00 =
\mathbf{8.62}$~atm.
(b) $T = PV/nR = (2.00)(4.00)/[(0.300)(0.08206)] = \mathbf{325}$~K.
(c) $n = 0.500/4.003 = \num{0.1249}$~mol;
$P = (0.1249)(0.08206)(298)/2.00 = \mathbf{1.53}$~atm.
(d) The pressure \textbf{doubles}. With $n$ and $V$ fixed, $P \propto T$
in kelvin, and the temperature doubled. In kinetic terms the molecules
move faster and strike the walls both more often and harder.\end{apcanswer}}
\end{yourturn}

% =====================================================================
\section*{\sffamily Ladder 6 \textbullet\ Gas stoichiometry}
\ced{3.4}

Moles are the bridge. Convert to moles, use the balanced equation, then
convert the moles of gas to a volume with $PV = nRT$. Use
\SI{22.4}{\litre\per\mole} \emph{only} at STP.

\begin{workedexample}[I do: solid to gas volume]
\SI{5.00}{\gram} of \ce{NaHCO3} ($M = \SI{84.01}{\gram\per\mole}$)
decomposes, releasing \ce{CO2}. What volume does the gas occupy at
\SI{1.00}{\atm} and \SI{298}{\kelvin}?

\textbf{Moles of solid:} $5.00/84.01 = \num{0.0595}$~mol.

\textbf{Through the equation:} each \ce{NaHCO3} releases one \ce{CO2}, so
$n(\ce{CO2}) = \num{0.0595}$~mol.

\textbf{To a volume:}
\[ V = \frac{nRT}{P} = \frac{(0.05952)(0.08206)(298)}{1.00}
   = \mathbf{1.46~L} \]

\textbf{Why not just use 22.4?} Because \SI{22.4}{\litre\per\mole} holds
only at STP (\SI{273}{\kelvin}, \SI{1}{\atm}). At \SI{298}{\kelvin} the
molar volume is \SI{24.5}{\litre}, and using 22.4 would give
\SI{1.33}{\litre} --- about \SI{9}{\percent} low.
\end{workedexample}

\begin{yourturn}[YOUR TURN 6 --- four questions]
\begin{enumerate}[leftmargin=*,itemsep=0.5em]
  \item What volume would that same \ce{CO2} occupy at STP?
        \workspace[1.6cm]
  \item \SI{0.200}{\mole} of \ce{O2} at \SI{2.00}{\atm} and
        \SI{273}{\kelvin}: volume? \workspace[1.6cm]
  \item Why is \SI{22.4}{\litre\per\mole} wrong at \SI{298}{\kelvin}?
        \workspace[1.6cm]
  \item Equal volumes of \ce{H2} and \ce{O2} at the same $T$ and $P$
        contain equal what? \workspace[1.6cm]
\end{enumerate}
\selfcheck{(a) \SI{1.33}{\litre} \quad (b) \SI{2.24}{\litre} \quad
(c) molar volume depends on $T$ and $P$ \quad (d) equal moles (equal
numbers of molecules)}
\keyonly{\begin{apcanswer}(a) At STP the molar volume is
\SI{22.4}{\litre\per\mole}: $0.05952 \times 22.4 = \mathbf{1.33}$~L.
(b) $V = nRT/P = (0.200)(0.08206)(273)/2.00 = \mathbf{2.24}$~L.
(c) Because molar volume is $RT/P$, which changes with both. At
\SI{298}{\kelvin} and \SI{1}{\atm} it is \SI{24.5}{\litre\per\mole}, not
\num{22.4}.
(d) Equal numbers of \textbf{moles}, hence equal numbers of molecules ---
Avogadro's principle. Note they do \emph{not} contain equal masses, since
the molar masses differ.\end{apcanswer}}
\end{yourturn}

% =====================================================================
\section*{\sffamily Ladder 7 \textbullet\ Partial pressures}
\ced{3.4}

Each gas in a mixture exerts the pressure it would exert alone. The total
is their sum, and each one's share is set by its \textbf{mole fraction} ---
not by its molar mass.

\begin{workedexample}[I do: mole fractions to partial pressures]
A vessel holds \SI{0.400}{\mole} \ce{N2}, \SI{0.200}{\mole} \ce{O2} and
\SI{0.400}{\mole} \ce{He} at a total pressure of \SI{3.00}{\atm}. Find
each partial pressure.

\textbf{Total moles:} $0.400 + 0.200 + 0.400 = \num{1.000}$.

\[ P_{\ce{N2}} = \frac{0.400}{1.000}(3.00) = \mathbf{1.20~atm} \qquad
   P_{\ce{O2}} = \frac{0.200}{1.000}(3.00) = \mathbf{0.60~atm} \]
\[ P_{\ce{He}} = \frac{0.400}{1.000}(3.00) = \mathbf{1.20~atm} \]

\textbf{Check:} $1.20 + 0.60 + 1.20 = 3.00$. \checkmark

\textbf{Note what did not matter:} helium is by far the lightest gas here
and nitrogen is seven times heavier, yet they exert identical partial
pressures because they are present in identical amounts. Pressure counts
particles, not mass.
\end{workedexample}

\begin{yourturn}[YOUR TURN 7 --- four questions]
\begin{enumerate}[leftmargin=*,itemsep=0.5em]
  \item A mixture is \SI{0.600}{\mole} \ce{Ar} and \SI{0.400}{\mole}
        \ce{Ne} at \SI{5.00}{\atm}. Partial pressure of \ce{Ar}?
        \workspace[1.6cm]
  \item Gas collected over water at \SI{25}{\celsius} reads
        \SI{760}{\torr} total. If water's vapour pressure is
        \SI{23.8}{\torr}, what is the dry-gas pressure?
        \workspace[1.6cm]
  \item Adding helium to a rigid vessel: what happens to the partial
        pressure of the gas already there? \workspace[1.6cm]
  \item Why does a heavier gas not exert a larger partial pressure at
        equal moles? \workspace[1.8cm]
\end{enumerate}
\selfcheck{(a) \SI{3.00}{\atm} \quad (b) \SI{736}{\torr} \quad
(c) unchanged \quad (d) pressure depends on particle count, not mass}
\keyonly{\begin{apcanswer}(a) $x(\ce{Ar}) = 0.600/1.000 = 0.600$;
$P = 0.600 \times 5.00 = \mathbf{3.00}$~atm.
(b) $760 - 23.8 = \mathbf{736}$~torr. Gas collected over water is
saturated with water vapour, which contributes its own partial pressure
and must be subtracted.
(c) \textbf{Unchanged.} Its partial pressure depends only on how many
moles of \emph{it} are present and the volume and temperature --- none of
which changed. Only the \emph{total} pressure rises.
(d) Because pressure arises from the momentum delivered by collisions with
the walls, and at a given temperature all gases have the same average
\emph{kinetic energy}. A heavier molecule moves more slowly in exactly the
compensating proportion, so equal numbers of particles exert equal
pressure.\end{apcanswer}}
\end{yourturn}

% =====================================================================
\section*{\sffamily Ladder 8 \textbullet\ Kinetic molecular theory and effusion}
\ced{3.5}

At a given temperature all gases have the same average \emph{kinetic
energy}. They do not have the same average \emph{speed} --- because
$\text{KE} = \tfrac{1}{2}mv^2$, the lighter molecule must move faster.

\begin{workedexample}[I do: energy, speed, and rate]
Helium and argon are at the same temperature. Compare their kinetic
energies and speeds, and find the ratio of their effusion rates.

\textbf{Kinetic energy: equal.} It depends only on absolute temperature.

\textbf{Speed: helium is faster.} With
$\tfrac{1}{2}mv^2$ the same for both, the smaller mass must carry the
larger speed.

\textbf{Effusion rate} goes as the inverse square root of molar mass:
\[ \frac{\text{rate}_{\ce{He}}}{\text{rate}_{\ce{Ar}}} =
   \sqrt{\frac{39.95}{4.003}} = \sqrt{9.98} = \mathbf{3.16} \]
Helium effuses about three times faster.

\textbf{Watch the direction of the ratio.} The \emph{lighter} gas is
faster, so the heavier mass goes on top. Getting it upside down gives
\num{0.32} --- a classic sign-of-the-error answer.
\end{workedexample}

\begin{yourturn}[YOUR TURN 8 --- four questions]
\begin{enumerate}[leftmargin=*,itemsep=0.5em]
  \item Ratio of effusion rates of \ce{H2} (2.016) to \ce{O2} (32.00):
        \workspace[1.6cm]
  \item Ratio for \ce{CH4} (16.04) to \ce{SO2} (64.07):
        \workspace[1.6cm]
  \item Two gases at the same temperature: which has the greater average
        kinetic energy? \workspace[1.4cm]
  \item Describe how the speed distribution of \ce{He} differs from that
        of \ce{Ar} at the same temperature. \workspace[2.0cm]
\end{enumerate}
\selfcheck{(a) \num{3.98} \quad (b) \num{2.00} \quad (c) equal \quad
(d) helium's peak is at higher speed, broader and lower}
\keyonly{\begin{apcanswer}(a) $\sqrt{32.00/2.016} = \mathbf{3.98}$.
(b) $\sqrt{64.07/16.04} = \sqrt{3.995} = \mathbf{2.00}$.
(c) \textbf{Equal} --- average kinetic energy depends only on absolute
temperature, not on identity.
(d) The helium curve \textbf{peaks at a much higher speed} and is
\textbf{broader and lower}, spreading its molecules over a wide range of
high speeds. Argon's peaks at a low speed and is narrower and taller. Both
enclose the \emph{same total area}, since both represent the same number
of molecules.\end{apcanswer}}
\end{yourturn}

% =====================================================================
\section*{\sffamily Ladder 9 \textbullet\ Deviation from ideal behaviour}
\ced{3.6}

The ideal model assumes molecules have no volume and no attractions. Both
assumptions fail --- and they fail in opposite directions.

\begin{workedexample}[I do: two regimes, two causes]
Explain why a real gas shows a lower pressure than predicted at moderately
high pressure, but a higher pressure at very high pressure.

\textbf{Moderately high pressure --- attractions dominate.} Molecules are
close enough to attract one another. One heading for the wall is pulled
back by those behind it and strikes with less force, so the measured
pressure falls \emph{below} the ideal prediction.

\textbf{Very high pressure --- molecular volume dominates.} Now the
molecules are packed so tightly that the space they themselves occupy is
no longer negligible. The free volume is less than $V$, so the gas is more
compressed than the law assumes and the pressure runs \emph{above} the
prediction.

\textbf{Where the model works best:} \textbf{high temperature} (kinetic
energy swamps the attractions) and \textbf{low pressure} (molecules far
apart, own volume negligible).
\end{workedexample}

\begin{yourturn}[YOUR TURN 9 --- four questions]
\begin{enumerate}[leftmargin=*,itemsep=0.5em]
  \item State the two conditions for most nearly ideal behaviour.
        \workspace[1.4cm]
  \item Which behaves more ideally at the same $T$ and $P$, \ce{He} or
        \ce{H2O} vapour? Why? \workspace[1.8cm]
  \item Which assumption fails when a gas liquefies on cooling?
        \workspace[1.6cm]
  \item Why does raising the temperature reduce the deviation caused by
        attractions? \workspace[1.8cm]
\end{enumerate}
\selfcheck{(a) high $T$, low $P$ \quad (b) \ce{He} --- no hydrogen bonding
\quad (c) the no-attractions assumption \quad (d) kinetic energy
overwhelms the attraction}
\keyonly{\begin{apcanswer}(a) \textbf{High temperature} and \textbf{low
pressure}.
(b) \textbf{Helium.} It is a tiny, nonpolar atom with only very weak
dispersion forces, so the no-attraction assumption nearly holds. Water
vapour hydrogen bonds strongly, so it deviates far more.
(c) The assumption of \textbf{no intermolecular attractions}. Liquefaction
happens precisely because attractions exist and become able to hold
molecules together once they slow down.
(d) At higher temperature the molecules have far more kinetic energy
relative to the strength of the intermolecular attractions. A passing
attraction barely alters a fast molecule's path, so the gas behaves as
though no forces acted.\end{apcanswer}}
\end{yourturn}

% =====================================================================
\section*{\sffamily Ladder 10 \textbullet\ Molarity and dilution}
\ced{3.7}

Molarity is moles of solute per litre of \emph{solution}, not of solvent.
On dilution the moles of solute are unchanged --- only the volume grows.

\begin{workedexample}[I do: a preparation and a dilution]
(i) \SI{0.0250}{\mole} of solute is made up to \SI{500.0}{\milli\litre}.
Find the molarity. (ii) Then \SI{50.0}{\milli\litre} of
\SI{2.00}{\Molar} stock is diluted to \SI{500.0}{\milli\litre}.

\textbf{(i)} $M = n/V = 0.0250/0.500 = \mathbf{0.0500~M}$.

\textbf{(ii)} Moles of solute do not change, so
$M_1V_1 = M_2V_2$:
\[ M_2 = \frac{(2.00)(50.0)}{500.0} = \mathbf{0.200~M} \]

\textbf{Why ``of solution'' matters:} to make the first solution you
dissolve the solid in \emph{less} than \SI{500}{\milli\litre} and then
make up to the mark. Adding solid to \SI{500}{\milli\litre} of water gives
a final volume above \SI{500}{\milli\litre} and a concentration below
target.
\end{workedexample}

\begin{yourturn}[YOUR TURN 10 --- four questions]
\begin{enumerate}[leftmargin=*,itemsep=0.5em]
  \item \SI{11.7}{\gram} \ce{NaCl} ($M = 58.44$) made up to
        \SI{250.0}{\milli\litre}. Molarity? \workspace[1.8cm]
  \item Moles of solute in \SI{40.0}{\milli\litre} of
        \SI{0.500}{\Molar}: \workspace[1.6cm]
  \item Dilute \SI{25.0}{\milli\litre} of \SI{6.00}{\Molar} to
        \SI{0.500}{\Molar}. Final volume? \workspace[1.6cm]
  \item After dilution, has the number of moles of solute changed?
        Explain. \workspace[1.6cm]
\end{enumerate}
\selfcheck{(a) \SI{0.801}{\Molar} \quad (b) \num{0.0200}~mol \quad
(c) \SI{300}{\milli\litre} \quad (d) no --- only the volume changed}
\keyonly{\begin{apcanswer}(a) $n = 11.7/58.44 = \num{0.2002}$~mol;
$M = 0.2002/0.2500 = \mathbf{0.801}$~M.
(b) $n = MV = (0.500)(0.0400) = \mathbf{0.0200}$~mol.
(c) $V_2 = M_1V_1/M_2 = (6.00)(25.0)/0.500 = \mathbf{300}$~mL.
(d) \textbf{No.} Adding solvent changes the volume, not the amount of
solute, so $n$ is constant and only the concentration falls. That
conservation is exactly what makes $M_1V_1 = M_2V_2$
valid.\end{apcanswer}}
\end{yourturn}

% =====================================================================
\section*{\sffamily Ladder 11 \textbullet\ Solubility and ``like dissolves like''}
\ced{3.10}

A solute dissolves when the solute--solvent attractions formed are
comparable to the solute--solute and solvent--solvent attractions broken.

\begin{workedexample}[I do: two solutes, one solvent]
Explain why \ce{NaCl} dissolves readily in water but not in hexane, and
why \ce{I2} does the reverse.

\textbf{\ce{NaCl} in water:} water is polar and surrounds each ion with
strong \textbf{ion--dipole} attractions --- oxygen ends toward \ce{Na+},
hydrogen ends toward \ce{Cl-}. The energy released repays the cost of
pulling the lattice apart.

\textbf{\ce{NaCl} in hexane:} hexane is nonpolar and offers only weak
dispersion forces, which cannot compensate for \ce{NaCl}'s very large
lattice energy. It does not dissolve.

\textbf{\ce{I2} reversed:} iodine is nonpolar. In hexane, dispersion
forces between solute and solvent are comparable to those being replaced,
so it dissolves. In water, dissolving would mean disrupting water's
hydrogen-bond network with nothing comparable offered in return --- so it
barely dissolves.
\end{workedexample}

\begin{yourturn}[YOUR TURN 11 --- four questions]
\begin{enumerate}[leftmargin=*,itemsep=0.5em]
  \item Will \ce{CH3OH} dissolve in water? Why? \workspace[1.6cm]
  \item Will \ce{CCl4} dissolve in water? Why? \workspace[1.6cm]
  \item Which is more soluble in water, \ce{CH3CH2OH} or
        \ce{CH3(CH2)6CH3}? \workspace[1.6cm]
  \item Describe how water molecules orient around a dissolved \ce{Ca^2+}
        ion. \workspace[1.8cm]
\end{enumerate}
\selfcheck{(a) yes --- hydrogen bonds with water \quad (b) no --- nonpolar
\quad (c) ethanol \quad (d) oxygen ends inward, toward the cation}
\keyonly{\begin{apcanswer}(a) \textbf{Yes.} Its \ce{-OH} group hydrogen
bonds with water, so the new solute--solvent attractions are comparable to
the water--water ones they replace.
(b) \textbf{No.} \ce{CCl4} is nonpolar and symmetric; it offers only
dispersion forces, which cannot compensate for disrupting water's
hydrogen-bond network.
(c) \textbf{Ethanol.} Both have a hydrocarbon portion, but ethanol's is
short and its \ce{-OH} dominates. Octane has no polar group at all.
(d) The water molecules turn their \textbf{oxygen} ends inward, toward the
cation. Oxygen carries the partial negative charge, so it is attracted to
the positive \ce{Ca^2+}; the attraction is \textbf{ion--dipole}, and it is
especially strong here because the ion carries a $2+$
charge.\end{apcanswer}}
\end{yourturn}

% =====================================================================
\section*{\sffamily Ladder 12 \textbullet\ Particulate pictures and separation}
\ced{3.8} \ced{3.9}

A particulate diagram of a solution must show what is actually in the
beaker: dissociated ions for a soluble salt, intact molecules for a
molecular solute, and only partial ionization for a weak acid.

\begin{workedexample}[I do: drawing three solutions correctly]
Describe correct particulate drawings of aqueous \ce{K2SO4},
\ce{C6H12O6} and \ce{HF}.

\textbf{\ce{K2SO4}} --- a soluble ionic compound, so it dissociates
completely. Draw \textbf{two} \ce{K+} for every one \ce{SO4^2-}, all
separated and hydrated. The sulfate stays intact as one unit: the
\ce{S-O} bonds are covalent and do not break.

\textbf{\ce{C6H12O6}} --- molecular and soluble. Draw whole, intact
glucose molecules dispersed among the water. \textbf{No ions at all.}

\textbf{\ce{HF}} --- a weak acid. Draw \emph{mostly} intact \ce{HF}
molecules with only a few \ce{H+} and \ce{F-} ions, reflecting partial
ionization.

\textbf{Separation follows the same logic:} filtration separates an
undissolved solid from a solution; distillation separates liquids by
volatility; chromatography separates by differing affinity for a
stationary phase.
\end{workedexample}

\begin{yourturn}[YOUR TURN 12 --- four questions]
\begin{enumerate}[leftmargin=*,itemsep=0.5em]
  \item How many particles result from one formula unit of \ce{CaCl2}
        dissolving? \workspace[1.4cm]
  \item Should a drawing of aqueous \ce{HCl} show intact \ce{HCl}
        molecules? Why? \workspace[1.6cm]
  \item Which technique separates sand from salt water?
        \workspace[1.4cm]
  \item Which technique separates ethanol from water, and on what
        property? \workspace[1.6cm]
\end{enumerate}
\selfcheck{(a) three --- one \ce{Ca^2+} and two \ce{Cl-} \quad (b) no ---
it is a strong acid, fully ionized \quad (c) filtration \quad
(d) distillation, by volatility}
\keyonly{\begin{apcanswer}(a) \textbf{Three}: one \ce{Ca^2+} and two
\ce{Cl-}.
(b) \textbf{No.} \ce{HCl} is a strong acid and ionizes essentially
completely, so the drawing should show separated \ce{H+} and \ce{Cl-} and
essentially no intact \ce{HCl}. (Contrast \ce{HF}, which is weak and
\emph{should} be drawn mostly intact.)
(c) \textbf{Filtration} --- the sand is undissolved and is retained by the
filter while the solution passes through.
(d) \textbf{Distillation}, exploiting the difference in
\textbf{volatility}: ethanol boils at \SI{78}{\celsius} and water at
\SI{100}{\celsius}, so the ethanol-rich vapour comes off
first.\end{apcanswer}}
\end{yourturn}

% =====================================================================
\section*{\sffamily Ladder 13 \textbullet\ Photons and the spectrum}
\ced{3.11} \ced{3.12}

$E = h\nu = hc/\lambda$. Energy is \emph{inversely} proportional to
wavelength: short wavelength means high energy.

\begin{workedexample}[I do: wavelength to energy per mole]
Find the energy of one photon of \SI{400}{\nano\metre} light, and the
energy per mole.

\[ E = \frac{hc}{\lambda} =
   \frac{(6.626\times10^{-34})(3.00\times10^{8})}{400\times10^{-9}}
   = \mathbf{4.97\times10^{-19}~J} \]

\textbf{Per mole:} $4.97\times10^{-19} \times 6.022\times10^{23} =
\num{2.99e5}$~J $= \mathbf{299~kJ/mol}$.

\textbf{The unit trap:} the wavelength must be in \textbf{metres}.
Leaving it in nanometres makes the answer wrong by a factor of
$10^{9}$.

\textbf{Sense check:} \SI{400}{\nano\metre} is violet, the high-energy end
of the visible range. Red light at \SI{700}{\nano\metre} should come out
lower --- and it does, at \SI{2.84e-19}{\joule}.
\end{workedexample}

\begin{yourturn}[YOUR TURN 13 --- four questions]
\begin{enumerate}[leftmargin=*,itemsep=0.5em]
  \item Energy of one photon of \SI{700}{\nano\metre} light:
        \workspace[1.8cm]
  \item Energy per mole of \SI{254}{\nano\metre} ultraviolet photons:
        \workspace[1.8cm]
  \item Which carries more energy per photon, microwave or ultraviolet?
        Why? \workspace[1.6cm]
  \item Doubling the intensity of a light source does not change whether
        an electron is ejected from a metal. Explain. \workspace[2.0cm]
\end{enumerate}
\selfcheck{(a) \SI{2.84e-19}{\joule} \quad (b) \SI{471}{\kilo\joule\per\mole}
\quad (c) ultraviolet --- shorter wavelength \quad (d) intensity is photon
\emph{number}, not photon energy}
\keyonly{\begin{apcanswer}(a) $E = (6.626\times10^{-34})
(3.00\times10^{8})/(700\times10^{-9}) = \mathbf{2.84\times10^{-19}}$~J.
(b) $E = \num{7.83e-19}$~J per photon;
$\times 6.022\times10^{23} = \mathbf{471}$~kJ/mol.
(c) \textbf{Ultraviolet}, by a very large factor. Energy goes as
$1/\lambda$, and UV wavelengths are thousands of times shorter than
microwave ones.
(d) Intensity sets how \emph{many} photons arrive per second, not how much
energy each one carries. Ejection depends on a \textbf{single} photon
having enough energy on its own to free an electron. More photons of the
same insufficient energy eject more of nothing.\end{apcanswer}}
\end{yourturn}

% =====================================================================
\section*{\sffamily Ladder 14 \textbullet\ The Beer--Lambert law}
\ced{3.13}

$A = \varepsilon b c$: absorbance is proportional to concentration at
fixed path length. That proportionality is what makes it an analytical
tool --- measure $A$, get $c$.

\begin{workedexample}[I do: absorbance to concentration]
A solution has $A = \num{0.620}$ in a \SI{1.00}{\centi\metre} cell, with
$\varepsilon = \SI{2.00e3}{}$~\si{\per\Molar\per\centi\metre}. Find the
concentration.

\[ c = \frac{A}{\varepsilon b} = \frac{0.620}{(2.00\times10^{3})(1.00)}
   = \mathbf{3.10\times10^{-4}~M} \]

\textbf{Why a calibration curve is used in practice:} you prepare
standards of known concentration, measure each one's absorbance, and plot
$A$ against $c$. The result is a straight line through the origin, and an
unknown's concentration is read off it. That avoids needing
$\varepsilon$ at all.

\textbf{Why absorbance and not transmittance:} absorbance is the quantity
that is \emph{linear} in concentration. Transmittance falls
exponentially, which is far harder to read.
\end{workedexample}

\begin{yourturn}[YOUR TURN 14 --- four questions]
\begin{enumerate}[leftmargin=*,itemsep=0.5em]
  \item Same solution diluted to half. New absorbance?
        \workspace[1.4cm]
  \item $c = \SI{1.5e-4}{\Molar}$, $\varepsilon = \SI{2.00e3}{}$,
        $b = \SI{1.00}{\centi\metre}$. Find $A$. \workspace[1.6cm]
  \item Why must the wavelength be chosen at the absorption maximum?
        \workspace[1.8cm]
  \item A student's cell is scratched and scatters light. What happens to
        the measured absorbance and the calculated concentration?
        \workspace[1.8cm]
\end{enumerate}
\selfcheck{(a) \num{0.310} \quad (b) \num{0.300} \quad (c) greatest
sensitivity \quad (d) both come out too high}
\keyonly{\begin{apcanswer}(a) \num{0.310}. $A \propto c$ at fixed path
length, so halving the concentration halves the absorbance.
(b) $A = (2.00\times10^{3})(1.00)(1.5\times10^{-4}) = \mathbf{0.300}$.
(c) At the absorption maximum, $\varepsilon$ is largest, so a given change
in concentration produces the largest change in absorbance --- the greatest
sensitivity. It is also the flattest part of the spectrum, so a small
error in wavelength setting matters least.
(d) Scattered light does not reach the detector, so the instrument reads
it as though it had been \emph{absorbed}: the measured absorbance is
\textbf{too high}, and the concentration calculated from it is
correspondingly \textbf{too high}.\end{apcanswer}}
\end{yourturn}

% =====================================================================
\section*{\sffamily Where you stand}

Tick a ladder only if all four YOUR TURN questions were right first time.

\renewcommand{\arraystretch}{1.30}
\noindent\begin{tabular}{@{}c p{6.0cm} c >{\raggedright\arraybackslash}p{4.9cm}@{}}
\toprule
\textbf{First try?} & \textbf{Skill} & \textbf{Ladder} &
  \textbf{If not, re-read\ldots}\\
\midrule
$\square$ & identifying IMFs & 1 & H must be on N, O or F \\
$\square$ & ranking boiling points & 2 & type first, then electrons \\
$\square$ & classifying solids & 3 & what breaks on melting \\
$\square$ & heating curves & 4 & flats are potential energy \\
$\square$ & ideal gas law & 5 & kelvin, always \\
$\square$ & gas stoichiometry & 6 & 22.4 L/mol is STP only \\
$\square$ & partial pressures & 7 & mole fraction, not mass \\
$\square$ & KMT and effusion & 8 & same $T$, same KE \\
$\square$ & real gases & 9 & attractions vs own volume \\
$\square$ & molarity and dilution & 10 & per litre of \emph{solution} \\
$\square$ & like dissolves like & 11 & compare forces made and broken \\
$\square$ & particulate pictures & 12 & dissociate only if it does \\
$\square$ & photons & 13 & $\lambda$ in metres \\
$\square$ & Beer--Lambert & 14 & $A \propto c$ \\
\bottomrule
\end{tabular}

\begin{apcnote}[Scoring yourself honestly]
14/14: outstanding --- this is the hardest unit in the course. 10--13:
solid; redo the missed ladders tomorrow from a blank page. 9 or fewer: the
recurring root causes here are exactly three --- (1) temperature left in
\si{\celsius}, (2) using \SI{22.4}{\litre\per\mole} away from STP, and
(3) asserting that one substance has ``stronger forces'' without naming
which force and why. Fix those three and most of these ladders fall
together.
\end{apcnote}

\end{document}
